\documentclass[../../../../main.tex]{subfiles}

\begin{document}

\section*{第1問}

\begin{proof}[解]
$a,b,p>0$とし，$t=\dfrac ba$（$t>0$）とおく。$C=A-B$とおくと
\[
C=a^p\left[(1+t)^p-2^{p-1}(1+t^p)\right]
\]
であり，この符号は$g(t)=(1+t)^p-2^{p-1}(1+t^p)$の符号に等しい。
\[
g'(t)=p(1+t)^{p-1}-p\cdot2^{p-1}t^{p-1}=p\left[(1+t)^{p-1}-(2t)^{p-1}\right]
\]

\textbf{$1^\circ$ $0<p<1$の時}

$p-1<0$だから$x\mapsto x^{p-1}$は（$x>0$で）単調減少。$0<t<1$では$1+t>2t$だから$(1+t)^{p-1}<(2t)^{p-1}$，よって$g'(t)<0$。$t>1$では$1+t<2t$だから$g'(t)>0$。$t=1$で$g'(1)=0$。
\begin{center}
\begin{tabular}{c|ccccc}
$t$ & $0$ & & $1$ & & \\
\hline
$g'$ & & $-$ & $0$ & $+$ & \\
\hline
$g$ & & $\searrow$ & $0$ & $\nearrow$ &
\end{tabular}
\end{center}
（$g(1)=2^p-2^{p-1}\cdot2=0$）より，$g(t)\ge0$（$t>0$）で等号は$t=1$のみ。したがって$C\ge0$，すなわち$A\ge B$，等号は$a=b$の時のみ。

\textbf{$2^\circ$ $p=1$の時}

$A,B$の表式から明らかに$A=B$。

\textbf{$3^\circ$ $1<p$の時}

$p-1>0$だから$x\mapsto x^{p-1}$は単調増加。$1^\circ$と逆に，$0<t<1$で$1+t>2t$より$(1+t)^{p-1}>(2t)^{p-1}$，$g'(t)>0$。$t>1$で$g'(t)<0$。
\begin{center}
\begin{tabular}{c|ccccc}
$t$ & $0$ & & $1$ & & \\
\hline
$g'$ & & $+$ & $0$ & $-$ & \\
\hline
$g$ & & $\nearrow$ & $0$ & $\searrow$ &
\end{tabular}
\end{center}
より$g(t)\le0$（$t>0$）で等号は$t=1$のみ。したがって$C\le0$，すなわち$A\le B$，等号は$a=b$の時のみ。

以上をまとめて，
\[
p=1\ \text{または}\ a=b\ \text{の時}\ A=B,\qquad
0<p<1\ \text{かつ}\ a\ne b\ \text{の時}\ A>B,\qquad
1<p\ \text{かつ}\ a\ne b\ \text{の時}\ A<B
\]
\end{proof}

\newpage

\section*{第2問}

\begin{proof}[解]
$O$から底面へ下ろした垂線の足を$H$（正$n$角形$A_1\cdots A_n$の中心）とし，$A_1H=x$（$0<x<1$）とおく。対称性から，
\[
V_n=n\times(\text{三角錐 } O\text{-}HA_1A_2) \quad\cdots\text{①}
\]
である。$\triangle OHA_1$は$H$が直角，斜辺$OA_1=1$だから，ピタゴラスの定理より
\[
OH=\sqrt{1-x^2} \quad\cdots\text{②}
\]
また，$\angle A_1HA_2=2\pi/n$（$H$を中心とした隣接2頂点のなす角）で，二等辺三角形$A_1A_2H$（$HA_1=HA_2=x$）の面積は
\[
\triangle A_1A_2H=\frac12x^2\sin\frac{2\pi}n=x^2\sin\frac\pi n\cos\frac\pi n \quad\cdots\text{③}
\]
$OH$は底面に垂直だから，三角錐$O$-$HA_1A_2$の高さは$OH$であり，②③から
\[
(O\text{-}HA_1A_2)=\frac13\sqrt{1-x^2}\cdot x^2\sin\frac\pi n\cos\frac\pi n \quad\cdots\text{④}
\]

①④から
\[
V_n=\frac n3\sin\frac\pi n\cos\frac\pi n\cdot x^2\sqrt{1-x^2}
\]

$x^2\sqrt{1-x^2}$を$x\in(0,1)$で最大化する。$t=x^2\in(0,1)$とおくと$x^2\sqrt{1-x^2}=\sqrt{t^2(1-t)}=\sqrt{t^2-t^3}$だから，$f(t)=t^2-t^3$の最大値を考えれば良い。$f'(t)=2t-3t^2=t(2-3t)$より，
\begin{center}
\begin{tabular}{c|ccccc}
$t$ & $0$ & & $2/3$ & & $1$ \\
\hline
$f'$ & & $+$ & $0$ & $-$ & \\
\hline
$f$ & & $\nearrow$ & & $\searrow$ &
\end{tabular}
\end{center}
より$t=2/3$（$x^2=2/3$）で最大。この時
\[
x^2\sqrt{1-x^2}=\sqrt{f(2/3)}=\frac23\sqrt{1-\frac23}=\frac23\cdot\frac1{\sqrt3}=\frac{2\sqrt3}9
\]
したがって，
\[
V_n=\frac n3\sin\frac\pi n\cos\frac\pi n\cdot\frac{2\sqrt3}9=\frac{2\sqrt3}{27}\,n\sin\frac\pi n\cos\frac\pi n
\]

\textbf{(2)} $n\sin(\pi/n)=\pi\cdot\dfrac{\sin(\pi/n)}{\pi/n}\to\pi$（$n\to\infty$），$\cos(\pi/n)\to1$だから，
\[
\lim_{n\to\infty}V_n=\frac{2\sqrt3}{27}\cdot\pi\cdot1=\frac{2\sqrt3\,\pi}{27}
\]
\end{proof}

\newpage

\section*{第3問}

\begin{proof}[解]
$\triangle ABC$を$AB=AC=1$，$BC=a$（$0<a<2$），$\angle CAB=\alpha$，$\angle ABC=\beta$とする。周上の2点$P,Q$を結ぶ線分が$\triangle ABC$の面積を2等分するとき，$P,Q$の位置は対称性から次の2通りに帰着できる。

\textbf{$1^\circ$ $P,Q$が辺$AB,AC$上にある場合}

$AP=x,\ AQ=y$（$0<x,y\le1$）とおく。$\triangle APQ$が$\triangle ABC$の面積の半分の時，
\[
xy=\frac12 \quad\cdots\text{①}
\]
$\triangle APQ,\ \triangle ABC$に余弦定理を用いて
\[
PQ^2=x^2+y^2-2xy\cos\alpha,\qquad a^2=1+1-2\cos\alpha
\]
後者から$\cos\alpha=\dfrac{2-a^2}2$，これと①を前者に代入して
\[
PQ^2=x^2+y^2-2xy\cdot\frac{2-a^2}2=x^2+y^2-\frac{2-a^2}2 \quad\cdots\text{②}
\]
（①より$2xy=1$を用いた。）①，$0<x,y\le1$のもとで$x^2+y^2$を最小化する。相加相乗平均から$x^2+y^2\ge2xy=1$，等号は$x=y$の時。①とあわせて$x=y=1/\sqrt2$（$\le1$をみたす）で最小値$x^2+y^2=1$をとる。②より
\[
\min PQ^2=1-\frac{2-a^2}2=\frac{a^2}2
\]

\textbf{$2^\circ$ $P,Q$が辺$AB,BC$上にある場合}（$AC,BC$上の場合も対称性から同じ結論になる）

$BP=x$（$0<x\le1$），$BQ=y$（$0<y\le a$）とおく。$\triangle BPQ$が半分の面積の時，$\triangle ABC$の面積が$\frac12\cdot1\cdot a\sin\beta$であることから
\[
xy=\frac a2 \quad\cdots\text{④}
\]
$\triangle BPQ,\ \triangle ABC$に余弦定理を用いて
\[
PQ^2=x^2+y^2-2xy\cos\beta,\qquad 1=1+a^2-2a\cos\beta
\]
後者から$\cos\beta=\dfrac a2$，④とあわせて
\[
PQ^2=x^2+y^2-2xy\cdot\frac a2=x^2+y^2-\frac{a^2}2 \quad\cdots\text{⑤}
\]
④，$0<x\le1,\ 0<y\le a$のもとで$x^2+y^2$を最小化する。相加相乗平均から$x^2+y^2\ge2xy=a$，等号は$x=y=\sqrt{a/2}$の時。この点が制約をみたすには$\sqrt{a/2}\le1$（$\iff a\le2$，常に成立）かつ$\sqrt{a/2}\le a$（$\iff a\ge1/2$）が必要。

\begin{itemize}
\item $\dfrac12\le a<2$の時：$x=y=\sqrt{a/2}$が制約内にあり，$\min(x^2+y^2)=a$。
\item $0<a<\dfrac12$の時：$\sqrt{a/2}>a$となり$y=\sqrt{a/2}$は$y\le a$に反するので，境界$y=a$（このもとで④から$x=1/2$）で最小となり，$\min(x^2+y^2)=\dfrac14+a^2$。
\end{itemize}
⑤とあわせて
\[
\min PQ^2=
\begin{cases}
\dfrac12a^2+\dfrac14 & \left(0<a\le\dfrac12\right)\\[2mm]
-\dfrac12a^2+a & \left(\dfrac12\le a<2\right)
\end{cases}
\]

\textbf{比較} $1^\circ$の$a^2/2$と$2^\circ$を比較する。$0<a\le1/2$では$2^\circ=a^2/2+1/4>a^2/2=1^\circ$。$1/2\le a<2$では
\[
2^\circ-1^\circ=\left(-\frac{a^2}2+a\right)-\frac{a^2}2=a(1-a)
\]
であり，$1/2\le a<1$で$\ge0$（$1^\circ$が小さい），$1<a<2$で$<0$（$2^\circ$が小さい），$a=1$で等しい。以上より，どちらの場合も$0<a\le1$では$1^\circ$（$PQ^2=a^2/2$）が，$1\le a<2$では$2^\circ$（$PQ^2=-a^2/2+a$）が全体の最小値を与える。したがって，もとめる線分の最小値は
\[
\min PQ=
\begin{cases}
\dfrac{\sqrt2}2\,a & (0<a\le1)\\[2mm]
\sqrt{-\dfrac12a^2+a} & (1\le a<2)
\end{cases}
\]
\end{proof}

\newpage

\section*{第4問}

$a_n=\displaystyle\int_0^1t^{2n-1}e^t\,dt$，$b_n=\displaystyle\sum_{k=1}^{n-1}\frac{{}_{2n-1}P_{2n-2k}}{2k+1}$とおき，
\[
a_n+e\,b_n=(2n-1)! \quad\cdots\Diamond
\]
を示す（$n\ge2$）。

\begin{proof}[解]
$n=2$の時，
\[
a_2=\int_0^1t^3e^t\,dt=\left[e^t(t^3-3t^2+6t-6)\right]_0^1=(-2e)-(-6)=6-2e
\]
\[
b_2=\frac13\,{}_3P_2=\frac13\cdot3\cdot2=2
\]
だから$a_2+eb_2=6-2e+2e=6=3!$となり，$\Diamond$は$n=2$で成立。

以下，$n=m\ (m\ge2)$で$\Diamond$の成立を仮定し，$n=m+1$での成立を示す。部分積分を2回用いて，
\[
a_{m+1}=\int_0^1t^{2m+1}e^t\,dt=\left[t^{2m+1}e^t\right]_0^1-(2m+1)\int_0^1t^{2m}e^t\,dt
\]
\[
=e-(2m+1)\left\{\left[t^{2m}e^t\right]_0^1-2m\int_0^1t^{2m-1}e^t\,dt\right\}=e-(2m+1)\{e-2m\,a_m\}
\]
\[
\therefore\ a_{m+1}=-2m\,e+(2m+1)(2m)a_m \quad\cdots\text{①}
\]

また，順列の定義${}_nP_r=n!/(n-r)!$から${}_{2m-1}P_{2m-2k}=(2m-1)!/(2k-1)!$であり，
\[
{}_{2m+1}P_{2m+2-2k}=\frac{(2m+1)!}{(2k-1)!}=(2m+1)(2m)\cdot\frac{(2m-1)!}{(2k-1)!}=(2m+1)(2m)\cdot{}_{2m-1}P_{2m-2k}
\]
だから
\[
b_{m+1}=\sum_{k=1}^m\frac{{}_{2m+1}P_{2m+2-2k}}{2k+1}=(2m+1)(2m)\sum_{k=1}^m\frac{{}_{2m-1}P_{2m-2k}}{2k+1}
\]
和を$k=m$の項（${}_{2m-1}P_0/(2m+1)=1/(2m+1)$）と残り（$=b_m$）に分けて，
\[
b_{m+1}=(2m+1)(2m)\left\{b_m+\frac1{2m+1}\right\}=2m+(2m+1)(2m)b_m \quad\cdots\text{②}
\]

①，②から
\[
a_{m+1}+e\,b_{m+1}=\{-2me+(2m+1)(2m)a_m\}+e\{2m+(2m+1)(2m)b_m\}
\]
\[
=(2m+1)(2m)(a_m+e\,b_m)=(2m+1)(2m)(2m-1)!\quad(\because\text{仮定})=(2m+1)!
\]
となり，$n=m+1$でも$\Diamond$が成立。以上，数学的帰納法により$\Diamond$は全ての$n\ge2$で成立する。
\end{proof}

\bigskip
\noindent\textbf{[解2]}

\begin{proof}[解2]
${}_{2n-1}P_{2n-2k}=(2n-1)!/(2k-1)!$だから
\[
b_n=(2n-1)!\sum_{k=1}^{n-1}\frac1{(2k+1)(2k-1)!}=(2n-1)!\sum_{k=1}^{n-1}\frac{2k}{(2k+1)!}
\]
（$(2k+1)!=(2k+1)(2k)(2k-1)!$より$2k/(2k+1)!=1/[(2k+1)(2k-1)!]$）。さらに$2k=(2k+1)-1$から
\[
b_n=(2n-1)!\sum_{k=1}^{n-1}\left\{\frac1{(2k)!}-\frac1{(2k+1)!}\right\}
\]
したがって，$\Diamond$を$(2n-1)!$で割った式として
\[
\Diamond\iff\frac{a_n}{(2n-1)!}+e\sum_{k=1}^{n-1}\left\{\frac1{(2k)!}-\frac1{(2k+1)!}\right\}=1 \quad\cdots\text{③}
\]
を示せば良い。①（$n=m$を単に$n$と書き直して）から
\[
\frac{a_{n+1}}{(2n+1)!}=\frac{a_n}{(2n-1)!}-\frac{2n}{(2n+1)!}\,e=\frac{a_n}{(2n-1)!}-e\left\{\frac1{(2n)!}-\frac1{(2n+1)!}\right\}
\]
（$2n/(2n+1)!=1/(2n)!-1/(2n+1)!$）。これを$n=1,2,\dots$と繰り返し用いて（$n\ge2$の時）
\[
\frac{a_n}{(2n-1)!}=a_1-e\sum_{k=1}^{n-1}\left\{\frac1{(2k)!}-\frac1{(2k+1)!}\right\}
\]
$a_1=\displaystyle\int_0^1te^t\,dt=\left[e^t(t-1)\right]_0^1=0-(-1)=1$だから，これは③に他ならない。以上により$\Diamond$が示された。
\end{proof}

\newpage

\section*{第5問}

\begin{proof}[解]
$e(\theta)=\cos\theta+i\sin\theta$とおく。

\textbf{(1)} $P_0=0,\ P_1=1+i$。$\overrightarrow{P_{n-1}P_n}$を表す複素数を$d_n$とすると，題意から
\[
d_{n+1}=\frac23e\left(\frac\pi3\right)d_n \quad\cdots\text{②}
\]
$d_1=1+i$（①）とあわせて
\[
d_n=\left\{\frac23e\left(\frac\pi3\right)\right\}^{n-1}(1+i)
\]
だから
\[
\overrightarrow{OP_n}=\overrightarrow{OP_0}+d_1+\cdots+d_n=\sum_{k=1}^n\left\{\frac23e\left(\frac\pi3\right)\right\}^{k-1}(1+i)=(1+i)\cdot\frac{1-\left\{\frac23e(\pi/3)\right\}^n}{1-\frac23e(\pi/3)}
\]
$\left|\dfrac23e(\pi/3)\right|=\dfrac23<1$だから$n\to\infty$で
\[
P_\infty=\frac{1+i}{1-\frac23e(\pi/3)}
\]
$e(\pi/3)=\dfrac12+\dfrac{\sqrt3}2i$より$\dfrac23e(\pi/3)=\dfrac13+\dfrac{\sqrt3}3i$，$1-\dfrac23e(\pi/3)=\dfrac23-\dfrac{\sqrt3}3i=\dfrac{2-\sqrt3i}3$だから
\[
P_\infty=\frac{3(1+i)}{2-\sqrt3i}=\frac{3(1+i)(2+\sqrt3i)}{(2-\sqrt3i)(2+\sqrt3i)}=\frac{3(1+i)(2+\sqrt3i)}7
\]
$(1+i)(2+\sqrt3i)=(2-\sqrt3)+(2+\sqrt3)i$だから
\[
P_\infty=\frac37\left\{(2-\sqrt3)+(2+\sqrt3)i\right\}
\]

\textbf{(2)} $Q_0=0,\ Q_1=z$。(1)と同様に，$\overrightarrow{Q_{n-1}Q_n}$を表す複素数を$\beta_n$とすると$\beta_1=z$，
\[
\beta_{n+1}=\frac12e\left(\frac\pi6\right)\beta_n \quad\therefore\ \beta_n=\left\{\frac12e\left(\frac\pi6\right)\right\}^{n-1}z
\]
\[
\overrightarrow{OQ_n}=\sum_{k=1}^n\beta_k=z\cdot\frac{1-\left\{\frac12e(\pi/6)\right\}^n}{1-\frac12e(\pi/6)} \ \xrightarrow{n\to\infty}\ Q_\infty=\frac z{1-\frac12e(\pi/6)}
\]
$e(\pi/6)=\dfrac{\sqrt3}2+\dfrac12i$より$\dfrac12e(\pi/6)=\dfrac{\sqrt3}4+\dfrac14i$，$1-\dfrac12e(\pi/6)=\dfrac{4-\sqrt3}4-\dfrac14i=\dfrac{(4-\sqrt3)-i}4$だから
\[
Q_\infty=\frac{4z}{(4-\sqrt3)-i}
\]

$Q_\infty=P_\infty$より
\[
\frac{4z}{(4-\sqrt3)-i}=\frac37\left\{(2-\sqrt3)+(2+\sqrt3)i\right\}
\]
\[
z=\frac3{28}\left\{(2-\sqrt3)+(2+\sqrt3)i\right\}\left\{(4-\sqrt3)-i\right\}
\]
右辺の$\{\ \}\{\ \}$を展開する（実部$=pr-qs$，虚部$=ps+qr$，ただし$p=2-\sqrt3,q=2+\sqrt3,r=4-\sqrt3,s=-1$として）：
\[
pr=(2-\sqrt3)(4-\sqrt3)=11-6\sqrt3,\qquad qs=-(2+\sqrt3)
\]
\[
pr-qs=11-6\sqrt3+2+\sqrt3=13-5\sqrt3
\]
\[
ps=-(2-\sqrt3),\qquad qr=(2+\sqrt3)(4-\sqrt3)=5+2\sqrt3
\]
\[
ps+qr=-(2-\sqrt3)+5+2\sqrt3=3+3\sqrt3=3(1+\sqrt3)
\]
よって
\[
z=\frac3{28}\left\{(13-5\sqrt3)+3(1+\sqrt3)i\right\}
\]
\end{proof}

\end{document}
